% §3 Index-Set Encoding Schemes

We now define the first parameterized map: the index-set scheme.

\subsection{Definition}

\begin{definition}[Index-set scheme]
\label{def:index-set-scheme}
An \emph{index-set encoding scheme} for $n$ modes is a triple of
functions $\scheme = (U, P, \mathrm{Occ})$:
\begin{align}
  U &: \{0, \ldots, n{-}1\} \times \{n\} \to 2^{\{0,\ldots,n{-}1\}},
    \label{eq:update-fn} \\
  P &: \{0, \ldots, n{-}1\} \to 2^{\{0,\ldots,n{-}1\}},
    \label{eq:parity-fn} \\
  \mathrm{Occ} &: \{0, \ldots, n{-}1\} \to 2^{\{0,\ldots,n{-}1\}}.
    \label{eq:occ-fn}
\end{align}
We call $U(j,n)$ the \emph{update set}, $P(j)$ the \emph{parity set},
and $\mathrm{Occ}(j)$ the \emph{occupation set} of mode~$j$.
\end{definition}

\subsection{The universal encoding algorithm}
\label{sec:universal-algorithm}

Given a scheme $\scheme$, the Majorana operators are constructed as
follows.  For each mode $j \in \{0, \ldots, n{-}1\}$, define two
Pauli strings of width~$n$:
\begin{align}
  S_j^{(c)} &= \bigotimes_{k=0}^{n-1}
    \begin{cases}
      X & \text{if } k \in \Occ{j} \setminus \Pset{j}, \\
      Y & \text{if } k \in \Occ{j} \cap \Pset{j}, \\
      Z & \text{if } k \in \Uset{j,n} \cup
          (\Pset{j} \setminus \Occ{j}), \\
      I & \text{otherwise},
    \end{cases}
    \label{eq:cj-from-scheme} \\[6pt]
  S_j^{(d)} &= \bigotimes_{k=0}^{n-1}
    \begin{cases}
      Y & \text{if } k \in \Occ{j} \setminus \Pset{j}, \\
      X & \text{if } k \in \Occ{j} \cap \Pset{j}, \\
      Z & \text{if } k \in \Uset{j,n} \cup
          (\Pset{j} \setminus \Occ{j}), \\
      I & \text{otherwise}.
    \end{cases}
    \label{eq:dj-from-scheme}
\end{align}
The ladder operators are then obtained
via~\eqref{eq:encoded-ladder}:
$\ad{j} = \tfrac{1}{2}(S_j^{(c)} - i S_j^{(d)})$.

\begin{remark}
Equations~\eqref{eq:cj-from-scheme}--\eqref{eq:dj-from-scheme}
constitute a single universal algorithm, implemented once.  All
variation between encodings is captured in the functions $U$, $P$,
$\mathrm{Occ}$.  The algorithm itself never changes.
\end{remark}

\subsection{Instantiation: Three classical encodings}

We now exhibit the three classical encodings as specific choices of
$(U, P, \mathrm{Occ})$.

\subsubsection{Jordan--Wigner}

\begin{equation}
  U_{\text{JW}}(j, n) = \emptyset, \qquad
  P_{\text{JW}}(j) = \{0, \ldots, j{-}1\}, \qquad
  \mathrm{Occ}_{\text{JW}}(j) = \{j\}.
  \label{eq:jw-scheme}
\end{equation}
The parity set consists of all preceding modes, producing the
characteristic $Z$-chain of Jordan--Wigner.  The update set is empty,
and the occupation set is a singleton.

\subsubsection{Parity}

\begin{equation}
  U_{\text{Par}}(j, n) = \{j{+}1, \ldots, n{-}1\}, \quad
  P_{\text{Par}}(j) =
    \begin{cases}
      \{j{-}1\} & j > 0, \\
      \emptyset & j = 0,
    \end{cases}
  \quad
  \mathrm{Occ}_{\text{Par}}(j) =
    \begin{cases}
      \{j{-}1, j\} & j > 0, \\
      \{0\} & j = 0.
    \end{cases}
  \label{eq:parity-scheme}
\end{equation}
The update set propagates changes to all higher qubits, while the
occupation set includes the previous mode.

\subsubsection{Bravyi--Kitaev}

The Bravyi--Kitaev encoding derives its index sets from the
\emph{Fenwick tree}~\cite{fenwick1994, seeley2012} on $n$ nodes:
\begin{align}
  U_{\text{BK}}(j, n) &= \text{ancestors of } j \text{ in the Fenwick tree}, \notag \\
  P_{\text{BK}}(j) &= \text{children of } j \text{ in the Fenwick tree}, \notag \\
  \mathrm{Occ}_{\text{BK}}(j) &= \text{descendants of } j \text{ (including } j\text{)}.
  \label{eq:bk-scheme}
\end{align}
The Fenwick tree structure ensures $|U_{\text{BK}}|, |P_{\text{BK}}|,
|\mathrm{Occ}_{\text{BK}}| \le \lceil \log_2 n \rceil$, yielding
$O(\log_2 n)$ Pauli weight.

\begin{proposition}[Encoding as data]
\label{prop:encoding-as-data}
Let $\mathfrak{S}_n$ denote the set of all index-set schemes for $n$
modes, and let $\mathfrak{E}_n$ denote the set of fermion-to-qubit
encodings satisfying CAR.  The universal
algorithm~\eqref{eq:cj-from-scheme}--\eqref{eq:dj-from-scheme} defines
a map $\encode : \mathfrak{S}_n \to \mathfrak{E}_n$.  This map is:
\begin{enumerate}
  \item[(a)] \textbf{Deterministic}: given $\scheme$, the encoding
    $\encode_\scheme$ is uniquely determined.
  \item[(b)] \textbf{Surjective over known encodings}: Jordan--Wigner,
    Bravyi--Kitaev, and Parity are all in the image of $\encode$.
  \item[(c)] \textbf{Compact}: each scheme is specified by at most
    $3n$ set values (three sets per mode).
\end{enumerate}
\end{proposition}

\begin{proof}
Part~(a) follows from the deterministic construction
in~\eqref{eq:cj-from-scheme}--\eqref{eq:dj-from-scheme}.  Part~(b) is
witnessed by~\eqref{eq:jw-scheme}--\eqref{eq:bk-scheme}.  Part~(c) is
immediate from the domain of the functions.
\end{proof}

\subsection{The SRL correctness conditions}

Not every triple of set functions produces an encoding satisfying CAR.
Seeley, Richard, and Love~\cite{seeley2012} identified conditions on
the sets that guarantee correctness:
\begin{enumerate}
  \item The sets $\Occ{j}$, $\Pset{j} \setminus \Occ{j}$, and
    $\Uset{j,n}$ are pairwise disjoint for each~$j$.
  \item For $i \neq j$: the qubit positions where $S_i^{(c)}$ and
    $S_j^{(c)}$ differ must have even overlap.
\end{enumerate}
Verifying these conditions for a candidate scheme is a finite check on
sets.  When the conditions hold, the resulting encoding satisfies CAR
(\cref{app:car-verification}).

\subsection{Structural limitation}

A natural question is whether one can \emph{derive} valid index sets
from a tree structure, making the index-set framework as general as the
tree framework of \cref{sec:tree-encoding}.  The SRL construction
attempts this by reading off update, parity, and occupation sets from
ancestors, children, and descendants of each node.  We prove in
\cref{sec:star-tree} that this derivation satisfies the CAR if and only
if the tree is a \emph{star} (depth~$\le 1$), a surprisingly
restrictive domain.  The Bravyi--Kitaev encoding survives via a
separate, Fenwick-specific construction.  This limitation motivates the
path-based framework, which is universal.

\begin{remark}[Defining an encoding]
\label{rem:defining-encoding}
In the index-set framework, defining a new encoding requires only
writing three set-valued functions and verifying the SRL conditions.
In a concrete implementation, this amounts to fewer than 10 lines of
code.  By contrast, existing software libraries require implementing
the entire mapping procedure from scratch for each encoding.
\end{remark}
